\section{Conclusion}
\label{sec:conclusion}

This study evaluated quantum molecular structure encoding (QMSE) --- using
BondOrderMatrix and CoulombMatrix representations embedded directly into
parameterised quantum circuits coupled to a TensorFlow dense head ---
for antimalarial activity prediction on the 19{,}849-molecule P3 benchmark.

\textbf{The stereo probe establishes a principal positive finding}: with
full-molecule encoding ($n_{\text{qubits}} = n_{\text{atoms}}$), QMSE
discriminates all five R/S stereoisomer pairs tested ($K(R,S) = 0.010$ vs.\
$K(R,R) = 1.000$; $\Delta K = 0.990$) and all three Z/E pairs
($K(E,Z) = 0.983$; $\Delta K = 0.017$) under exact Aer statevector
simulation. This is the first direct kernel-level verification that
BondOrderMatrix encoding encodes chirality and geometric isomerism into
distinguishable quantum states. The Kernel Target Alignment ($= 0.242$)
on the P1 Set~A 17-molecule panel confirms non-trivial label alignment
of the QMSE kernel matrix.

\textbf{The canonical Phase~D benchmark yields a Scenario~C result}: the
QMSE--TF hybrid (CoulombMatrix, 8~qubits, depth~2, IsingYY, $h = 128 \times 3$,
Optuna best configuration) achieves AUC $= 0.640 \pm 0.044$ under canonical
5-fold CV on the frozen $n = 1{,}000$ P3 subsample, $\Delta = -0.229$
below the matched ECFP4-RBF SVM ($0.869 \pm 0.027$; $p = 0.0009$).
This underperformance is mechanistically attributable to the qubit-count
ceiling ($n_{\text{qubits}} = 8$), which retains fewer than half the heavy
atoms of typical ANP scaffolds and places stereocenters at indices $\geq 8$
outside the encoding window. The expressibility--trainability trade-off at
8 qubits and the limited training budget ($\approx 800$ samples, 80 epochs)
further constrain convergence.

An archived ECFP4-VQC result ($0.847 \pm 0.013$, Phase~4) provides a
Scenario~B (competitive, not superior) baseline for the classical-feature
encoding estimand. The comparison table is:
ECFP4-RBF (P3 canonical) $0.9475$;
Hybrid RF $0.8876$;
ECFP4-VQC ($0.8474$, archived EXPLORATORY);
QKS $0.8385$;
QMSE--TF hybrid $0.640$ (CANONICAL).

Taken together, the P3/P5/P7 series now documents three distinct
representation strategies --- quantum-inspired kernels on ECFP4 inputs (P3,
Scenario~B), GNNs/transformers on scaffold split (P5, Scenario~C),
and structure-direct QMSE hybrid circuits (P7, Scenario~C) --- all
yielding honest-negative results at the current NISQ/statevector scale.
The common bottleneck is the qubit-count ceiling, not the encoding
principle: QMSE \emph{does} discriminate stereoisomers in the quantum
kernel; the TF hybrid cannot leverage this because the $n_{\text{qubits}} = 8$
truncation erases most stereocenter information for drug-sized ANP molecules.

\textbf{Recommended next steps.}
The two complementary paths for overcoming the qubit ceiling are:
(i) quantum kernel SVM with Nystr\"{o}m approximation on IBM Quantum hardware,
using $n_{\text{qubits}} = n_{\text{atoms}}$ (full-molecule encoding, no truncation)
on a landmark set of $\approx 500$ P3 molecules --- this directly exploits the
stereo-discriminating capability confirmed here at the kernel level; and
(ii) atom-selection strategy for the TF hybrid, encoding the
$n_{\text{qubits}} = 8$ atoms adjacent to the stereocenter rather than the
first 8 atoms, which would restore stereochemical signal for the majority
of ANP scaffolds without requiring additional qubits.

Hardware validation on IBM Quantum devices remains an open task. The IBM
token and circuit infrastructure are in place; access depends on hardware
allocation. The stereo-probe circuits (circuit depth 22 at $n_q = 17$) and
the P1 Set~A kernel matrix are validated and ready for hardware submission.
